\chapter{Introduction and Background} \label{Chap2}
\pagenumbering{arabic}

\section{Introduction}

Deep reinforcement learning has achieved remarkable progress in robot control, yet its success depends heavily on the quality of the reward function. In practice, designing a reward function that both guides learning and is fully aligned with the true objective is extremely difficult\cite{ibrahim2024comprehensive}: in RL-based fine-tuning, the reward signal tends to be either too sparse, giving feedback only upon task completion\cite{andrychowicz2017hindsight}---especially in long-horizon manipulation tasks, where intermediate steps lack explicit reward guidance\cite{wang2022learning,yang2026actor}---or not fully aligned with the true task objective, containing misleading shaping components\cite{karwowski2024goodhart}. The common feature of these two cases is that the reward signal cannot reliably point to correct behaviour: even when the reward is dense in value, the inconsistency between a shaping reward and the true objective is itself a form of unreliable information; a high reward does not imply that the task is progressing in the right direction.

Meanwhile, pretraining--fine-tuning has become an important paradigm for building robotic policies\cite{firoozi2025foundation}, in which a general policy is first pretrained on large-scale data and then fine-tuned on a target task to adapt it to a specific embodiment and task; RL-based fine-tuning is a key component of this process. PPO\cite{schulman2017proximal} is currently the most widely used on-policy RL algorithm and is often the default choice for fine-tuning. Its advantage estimation relies on GAE\cite{schulman2015high}: if a shaping component consistently increases the return, then even if it conflicts with the final task objective, it can exert persistent optimization pressure. Consequently, when the reward condition during fine-tuning deviates from the true objective, PPO's optimization direction may be pulled by misleading reward components.

In recent years, GRPO has been widely adopted and achieved significant success as an alternative in RL fine-tuning of LLMs\cite{shao2024deepseekmath}. Unlike PPO, which takes raw reward maximization as its optimization basis, GRPO changes the way rewards are compared, shifting from absolute improvement to within-group relative improvement: when a reward component is generally increased across a group of samples, within-group normalization weakens its role as a discriminating signal, so GRPO may become less sensitive to systematic reward inflation\cite{yao2026multimodal}. In robot fine-tuning, an increased reward does not necessarily imply task progress, and this difference may make GRPO more insensitive to systematic reward growth that is inconsistent with the task objective. However, the application of GRPO in robotics is still in its infancy, and a systematic comparison with PPO under unreliable reward conditions in fine-tuning is still lacking.

To fill this gap, this thesis systematically compares the fine-tuning behaviour of PPO and GRPO on the Door task of HumanoidBench\cite{sferrazza2024humanoidbench} under two reward conditions: a Misaligned Reward condition that simulates unreliable reward information, and a Sparse Reward condition that simulates extremely sparse rewards. We find that PPO undergoes systematic reward hacking when rewards are unreliable, and collapses or stalls when rewards are sparse; GRPO remains stable under both conditions, with its policy-update magnitude varying with seed and condition while its return stays stable. Diagnostic analysis shows that this stability stems from the normalization effect of group-relative advantage; KL-penalty ablation and cross-task validation further confirm the source and generality of this mechanism. Our main contributions are:

\begin{enumerate}
  \item We reveal the empirical reward-hacking behaviour of PPO under unreliable rewards. Under the Misaligned condition PPO's return soars from a baseline of 244 to 431 while the success rate drops from 5.0\% to 0.0\%; under the Sparse condition it either collapses or stalls, with the failure mode switching with the reward condition.
  \item We demonstrate the stability of GRPO under both reward conditions and its KL-return decoupling. GRPO's return is stable within a narrow band of 188--219 across all seeds; its update magnitude varies with seed and reward condition, yet its return is decoupled from KL and does not drift with the update magnitude.
  \item Through reward decomposition and KL-penalty ablation we diagnose the source of the stability mechanism. Component-wise reward decomposition shows that PPO's return gain is concentrated in components weakly related to success; the KL-penalty ablation rules out the alternative explanation that ``stability stems from the absence of the KL penalty,'' attributing stability to the normalization effect of group-relative advantage.
  \item We provide cross-task evidence and practical implications. Additional experiments on the Reach task show that the conclusions are not limited to the Door task; in scenarios where reward design is imperfect or rewards are extremely sparse, GRPO is a more robust fine-tuning choice than PPO, and model selection should jointly monitor reward component decomposition and success rate rather than return increases alone.
\end{enumerate}

\section{Related Work}

\subsection{Misaligned Rewards and Reward Hacking}

The goal of reinforcement learning is set by the reward function. However, in real tasks the reward function often cannot perfectly describe the true objective and can only serve as an imperfect proxy for it; over-optimizing such an imperfect proxy ends up harming performance on the true objective\cite{karwowski2024goodhart}. When the reward design is not fully aligned with the true objective, an agent may systematically exploit loopholes in the reward function and optimize a surrogate metric instead of the true objective; this phenomenon is called reward hacking\cite{amodei2016concrete}. Skalse et al.\cite{skalse2022defining} provide a formal definition: optimizing an imperfect proxy causes a drop in performance under the true reward. Everitt et al.\cite{everitt2017reinforcement} were the first to formalize this problem as a Corrupt Reward MDP, pointing out that reward signals may be corrupted by, e.g., sensor errors, and that traditional RL methods struggle to compensate. Pan et al.\cite{pan2022effects} show in a systematic study across four misspecified-reward environments that stronger agents are more prone to exploit misspecified rewards, and that a phase transition exists in which true rewards drop sharply once capability exceeds a threshold. As model scale grows, reward hacking is equally prevalent in LLMs and multimodal models; the survey by Wang et al.\cite{wang2026reward} unifies these phenomena under the view of high-dimensional objectives being compressed into reward representations. In continuous control and robot manipulation tasks, every component of reward shaping is a potential target of exploitation, making the evaluation and mitigation of reward hacking an important problem for safe reinforcement learning.

\subsection{Sparse Rewards and Long-Horizon Tasks}

When the reward signal is provided only upon final task completion, the agent receives no learning signal on the vast majority of time steps; this setting is called a sparse reward\cite{andrychowicz2017hindsight}. It introduces two kinds of difficulty: exploration efficiency is extremely low, since random exploration can rarely reach a successful state\cite{wang2022learning}; and credit assignment is hard, since even on accidental success it is difficult to tell which early decisions led to the success\cite{yang2026actor}. When the task further has a long-horizon, multi-stage structure, both difficulties are amplified\cite{wang2022learning,yang2026actor}. The robotics community has therefore developed a variety of methods to recover learning signals, such as hindsight experience replay\cite{andrychowicz2017hindsight}, guidance from demonstrations\cite{vecerik2017leveraging}, reuse of off-the-shelf base controllers\cite{wang2022learning}, action chunking with intrinsic rewards\cite{yang2026actor}, curriculum learning\cite{xie2020learning}, curiosity-driven exploration\cite{hare2019dealing,schwarke2023curiosity}, and hierarchical structure\cite{hansen2025hierarchical}. In contrast, our Sparse Reward setting deliberately avoids such additional guidance and directly investigates whether the choice of the policy-optimization algorithm itself affects the robustness of a policy to extremely sparse feedback.

\subsection{Pretraining--Fine-Tuning and Robot RL Post-Training}

In this paradigm, RL post-training is a key tool in the fine-tuning stage. On the one hand, RL has been used to train vision--language--action foundation models: SimpleVLA-RL\cite{li2026simplevla} scales the training of a VLA model via online reinforcement learning, and ReinFlow\cite{zhang2026reinflow} fine-tunes a flow-matching policy with online RL. On the other hand, purely RL-pretrained policies also rely on fine-tuning to adapt to new environments, e.g., large-scale SAC pretraining followed by model-based fine-tuning\cite{huang2026towards} and transformer sequence-modeling pretraining followed by RL fine-tuning over uneven terrain\cite{radosavovic2024learning}. Robot Trains Robot\cite{hu2025robot} further studies fine-tuning in the real world, where an arm teacher assists a humanoid student. The survey by Gu et al.\cite{gu2026humanoid} systematically reviews these learning-based methods in humanoid locomotion and manipulation. These works are usually evaluated on their own collected data and benchmarks, and a unified humanoid environment for systematically studying behaviour in the fine-tuning stage is lacking. HumanoidBench\cite{sferrazza2024humanoidbench} provides a standardized evaluation environment for this paradigm; the Door task used in this thesis comes from this benchmark. In these paradigms, fine-tuning a pretrained policy on a target environment is common practice, but the reward condition during fine-tuning may deviate from that during pretraining, and the effect of such a deviation on policy stability has not been systematically studied.

\subsection{GRPO and Its Behaviour under Unreliable Rewards}

GRPO was originally proposed for mathematical reasoning in LLMs\cite{shao2024deepseekmath}; its group-based comparison mechanism is naturally robust to individual-level reward noise. Noise-corrected GRPO\cite{mansouri2025noise} analyzes this property theoretically and proposes a noise-correction method to obtain unbiased gradients; Yao et al.\cite{yao2026multimodal} observe in multimodal RL alignment experiments that GRPO is the most reward-hacking-resistant among three compared algorithms. In continuous control, Khanda et al.\cite{khanda2025extending} present the first theoretical framework extending GRPO to continuous control, with components including trajectory-clustering grouping and state-aware advantage estimation. De Oliveira et al.\cite{de2025learning} observe in classical RL environments that, on long-horizon continuous-control tasks, critic-free GRPO is usually inferior to PPO, with exceptions only in short-horizon environments, indicating that GRPO's advantage is not universal. Doorman\cite{xue2026opening} applies GRPO to sim-to-real fine-tuning of a humanoid Door task to mitigate partial observability under binary success signals. Taken together, existing work mainly focuses on whether GRPO can be used for continuous control or specific robot scenarios, or on comparing GRPO and PPO when used directly for training from scratch; although some work notices the application of GRPO in fine-tuning scenarios, a systematic study of the relative behaviour of GRPO and PPO under unreliable and extremely sparse rewards is still lacking. This thesis addresses this gap by systematically comparing the relative behaviour of the two algorithms in a fine-tuning setting, and further analyzing the relationship between group-relative advantage and the stability of policy updates.
